\documentclass[tc, manuscript]{copernicus}

\graphicspath{{../}}

\begin{document}

\title{A temperature-free morphological observable that distinguishes non-local from down-gradient heat transport at a melting interface}

% Corresponding-author email: author's AIET institutional account.
\Author[1]{Abdelrahman}{Saifelden}

\affil[1]{Alexandria Higher Institute of Engineering and Technology (AIET), Alexandria, Egypt}

\runningtitle{A temperature-free morphological test of interface heat transport}
\runningauthor{A. Saifelden}
\correspondence{Abdelrahman Saifelden (Abdulrahman.Saifelden22011411@aiet.edu.eg)}

\received{}
\pubdiscuss{}
\revised{}
\accepted{}
\published{}

\firstpage{1}
\maketitle

\begin{abstract}
Ice and rock scallops migrate downstream as they evolve, a motion classically tied to the phase lag between the wall heat flux and the interface topography: a melting surface migrates downstream when the flux maximum lies downstream of the trough, i.e. when the flux–topography phase lies in $(\pi /2, \pi )$ (Ashton \& Kennedy 1972; Gilpin et al. 1980; Hanratty 1981). A local, memoryless, down-gradient (``K-theory'') closure carries no such lag. Freezing a sinusoidal interface $y(x)=\bar{y}+a \sin(Kx)$ and projecting the time-averaged wall flux onto in-phase (amplitude rate $\beta $) and quadrature (migration $\omega _{mig}$) harmonics, every coefficient is solver-measured: the conduction amplitude rate is wavelength-independent ($\beta /a \sim  K^{+0.13}$), falsifying the $K^{2}$ and Mullins–Sekerka $|k|$ ansätze; the in-phase excess is negative at every amplitude and drive (smoothing-limited); and the quadrature grows sub-kinematically ($E_{cos} \propto  U^{\approx 0.5}$, friction-velocity controlled), giving a damped, downstream-migrating mode $s = -\beta  + i\omega _{mig}$. Swapping only the heat-transport operator at fixed momentum drives the quadrature to machine zero ($-1.2\times 10^{-18}$ for a uniform eddy diffusivity) against $-1.1\times 10^{-4}$ for resolved advection. We then construct a temperature-free observable $I \equiv  \tau \cdot c_{mig}/\lambda  = |Im(s)|/(2\pi |Re(s)|)$ in which conductivity, latent heat, and basal temperature drop cancel, computable from interface morphology alone; the solver pins $I = O(0.3\mbox{--}0.9)$, and a regime-matched reading of Bushuk et al. (2019) gives $I_{obs} \approx  0.05\mbox{--}0.4$ with the predicted downstream sign. We give a field-test protocol and a self-validated reanalysis harness (recovers $I$ to $<3\%$); the field test is proposed, not performed here, as the raw $h(x,t)$ arrays are not openly available.
\end{abstract}

\introduction

Scalloped erosion/melt patterns are ubiquitous on soluble and meltable surfaces — cave walls, glacier soffits, ice-shelf bases, meteorites, and lab ice (Curl 1966; Blumberg \& Curl 1974; Ashton \& Kennedy 1972; Bushuk et al. 2019). Their wavelength is fluid-selected through a friction-Reynolds criterion $Re_*=u_*\lambda /\nu \approx 2200$. A robust but under-exploited observation is that the pattern \textbf{migrates downstream} as it evolves — maximum ablation sits downstream of the trough, at flow reattachment (Blumberg \& Curl 1974; Gilpin, Hirata \& Cheng 1980).

The migration direction and its phase-lag origin are established (Gilpin et al. 1980; Hanratty 1981): downstream migration corresponds to a flux–topography phase in $(\pi /2,\pi )$, and a local, memoryless, down-gradient closure carries no such phase. What is \textbf{new here} is not that fact but three things built on it: a phase-separated harmonic decomposition that returns the amplitude and migration rates as independent solver measurements; a direct operator-swap experiment that quantifies how completely a down-gradient closure discards the migration-carrying asymmetry; and a temperature-free morphological observable that makes the phase lag measurable from field geometry alone, without knowing the basal temperature drop or any material constant. The non-local, direction-selecting structure the migration requires is exactly the structure (memory, backscatter) that a local, down-gradient closure cannot represent.

\textbf{Contributions.}

\begin{enumerate}
\item A \textbf{phase-separated} wall-flux harmonic decomposition that returns the corrected mode $s=-\beta +i\omega _{mig}$ as an \emph{output}, not an ansatz, with the conduction amplitude rate measured wavelength-independent — falsifying the $K^{2}$/$|k|$ curvature-smoothing ansätze for the conduction baseline, generalised across amplitude and drive (§3–§4).
\item A \textbf{direct operator-swap measurement} quantifying how far down-gradient closures fall short of the known phase-lag requirement: swapping only the heat-transport operator collapses the migration signal to machine zero (§5).
\item A \textbf{temperature-free, constant-free morphological observable} $I=\tau \cdot c_{mig}/\lambda $, with an implemented, self-validated reanalysis harness and a regime-matched literature bound, that makes the phase lag measurable from field morphology alone (§6–§8).
\end{enumerate}

\section{Method — harmonic decomposition of the wall flux}

Freeze $y(x)=\bar{y}+a \sin(Kx)$ ($K=2\pi n_{w}/L_{x}$), time-average the per-column melt flux $m(x)=-\kappa  \partial \theta /\partial y$, and project onto the two corrugation harmonics:

\begin{quote}
$E_{sin} = 2\langle e(x) \sin(Kx)\rangle $ (in-phase with the \textbf{shape} \ensuremath{\to} amplitude change, $\beta $) $E_{cos} = 2\langle e(x) \cos(Kx)\rangle $ (in \textbf{quadrature} \ensuremath{\to} pattern \textbf{migration}, $\omega _{mig}$)
\end{quote}

$E_{cos}$ is the Curl-1966 reattachment/lee signature: zero by symmetry unless the flow carries a direction. Two decompositions are run: the conduction baseline $m_{cond}$ (flow off) for the $\beta  \sim  K^{p}$ scaling, and the flow-induced excess $e = m_{flow} - m_{cond}$ versus drive $U$. Freezing the interface removes boundary motion so the flux is uncontaminated; this is \emph{why} the two coefficients separate cleanly, where a moving-boundary fit conflates them and returns unphysical (both-negative) coefficients.

\section{Findings}

\begin{center}\small
\begin{tabularx}{\linewidth}{>{\raggedright\arraybackslash}X>{\raggedright\arraybackslash}X>{\raggedright\arraybackslash}X}
\hline
quantity & measured & ansatz it kills \\
\hline
conduction $\beta /a \sim  K^{+0.13}$ & wavelength-\textbf{independent} & $K^{2}$ curvature and Mullins–Sekerka $|k|$ \textbf{falsified} \\
in-phase flow excess $E_{sin}$ & \textbf{negative at every driven $(a,U)$} & no autonomous $+\alpha  a^{1/2}$ growth — \textbf{smoothing-limited} \\
quadrature $E_{cos}$ & $\approx 0$ at $U=0$, grows $\propto  U^{0.5\mbox{--}0.8}$ & migration is real and \textbf{friction-velocity} ($u_*\sim U^{1/2}$) controlled, not kinematic ($\propto U$) \\
\hline
\end{tabularx}
\end{center}

\textbf{Corrected mode:} $s(K,U) = -\beta (K,U) + i\cdot \omega _{mig}(U)$, $Re(s)<0$ always, $Im(s)\propto U^{1/2}$ — a damped, downstream-migrating pattern, not a growth–saturation balance.

\textbf{Why this counts as a result, not a check.} The corrected mode \emph{fell out of} the decomposition; every coefficient is solver-measured; physical units are pinned independently by the Curl selection law without choosing $\lambda $ by hand; and the final ratio is constant-free and falsifiable.

\section{Generalisation across amplitude and drive}

\textbf{4.1 Amplitude ($a_{0}/\lambda  \in  [0.05, 0.40]$).} The two structural verdicts are not amplitude artefacts: the driven in-phase excess is \textbf{negative at every amplitude} (smoothing throughout, magnitude growing monotonically), and the conduction exponent stays $p\in [-0.27,+0.66]$, far below the $K^{2}$ ($+2$) and $|k|$ ($+1$) ansätze at every amplitude. \emph{Honest caveat:} strict $K$-independence ($|p|<0.6$) and amplitude-flat $\beta /a$ hold only in the signal-rich regime $a_{0}/\lambda \gtrsim 0.2$; at shallow amplitude the single-wavenumber conduction signal is at the noise floor (the K\ensuremath{^{2}}/|k| falsification is robust regardless). (\texttt{amplitude\_\allowbreak{}generalization\_\allowbreak{}scan}.)

\textbf{4.2 Drive ($U$ to 6, beyond the swept $[1.5,3.0]$).} The genuine falsification risk is that strong-drive lee separation adds an \emph{in-phase} component (an emergent growth term the weak-drive sweep would miss). It does not: $E_{sin}$ is \textbf{negative at every $U$} out to $U=6$, and the migration follows a \textbf{sub-kinematic} $E_{cos} \propto  U^{+0.48}$ (\ensuremath{\approx}½), far below the kinematic-advection alternative $U^{1}$. \emph{Honest caveat:} migration is non-monotone — it peaks near $U\approx 4.5$ and rolls over by $U=6$ (the lee structure reaches a limiting form), which \emph{reinforces} the sub-kinematic reading. (\texttt{drive\_\allowbreak{}window\_\allowbreak{}scan}.)

\begin{center}\small
\begin{tabular}{lll}
\hline
$U$ & in-phase $E_{sin}$ & quadrature $E_{cos}$ \\
\hline
0.0 & $-5.17\times 10^{-6}$ (smoothing) & $+4.9\times 10^{-8}$ (\ensuremath{\approx}0, parity control) \\
1.5 & $-6.23\times 10^{-5}$ & $-6.59\times 10^{-5}$ \\
3.0 & $-9.19\times 10^{-5}$ & $-1.14\times 10^{-4}$ \\
4.5 & $-8.86\times 10^{-5}$ & $-1.38\times 10^{-4}$ \\
6.0 & $-4.81\times 10^{-5}$ & $-1.19\times 10^{-4}$ \\
\hline
\end{tabular}
\end{center}

\section{A direct operator-swap measurement of the closure shortfall}

A local, memoryless, down-gradient eddy diffusivity acting on a sinusoid is symmetric under $x\to -x$: it carries no flow direction and can only return an in-phase (real, damping) response — $E_{cos}\equiv 0$. This is the closure-space statement of the classical phase-lag requirement (Gilpin et al. 1980; Hanratty 1981). The resolved advective flux measured $E_{cos}\ne 0$, so migration requires the flow-direction-selecting, non-local reattachment asymmetry that K-theory discards. The contribution here is to make the shortfall \textbf{quantitative}: advancing the \textbf{same} momentum field at the \textbf{same} drive ($U=3$, same interface) but transporting heat with the K-theory flux instead of resolved advection measures exactly how far each closure falls below the resolved migration:

\begin{center}\small
\begin{tabularx}{\linewidth}{>{\raggedright\arraybackslash}X>{\raggedright\arraybackslash}X>{\raggedright\arraybackslash}X}
\hline
heat transport (same $U=3$, same interface) & $E_{cos}$ & reading \\
\hline
resolved advective $-u\cdot \nabla \theta $ & $-1.14\times 10^{-4}$ & migrates \\
\textbf{uniform} eddy diffusivity $\nabla \cdot (K\nabla \theta )$ & $-1.2\times 10^{-18}$ & machine zero ($\sim 10^{-14}$ of resolved, $std=0$) \\
\textbf{flow-aware} Smagorinsky $K_{eddy}(|S|)$ & $+2.0\times 10^{-6}$ & $\sim 57\times $ below resolved \\
\hline
\end{tabularx}
\end{center}

Both closures still \textbf{smooth} ($E_{sin}\ne 0$): K-theory can damp but cannot migrate. GPU/CuPy controls (Tesla P100) confirm (a) parity — $E_{cos}\to +2.1\times 10^{-6}$ at $U=0$ vs $-8.4\times 10^{-5}$ at $U=3$ (\textasciitilde{}40\ensuremath{\times} smaller); and (b) the high-$Re_*$ migration exponent tightens to \textbf{$0.52$} ($nx=160$, four seeds, scatter \textasciitilde{}1–2 \%), on the friction-velocity $\tfrac{1}{2}$. (\texttt{scallop\_\allowbreak{}ktheory\_\allowbreak{}control.\allowbreak{}py}.)

This is the morphological analogue of the memory and backscatter signatures that a down-gradient closure omits in \emph{flux} space: the closure's blindness does not merely lose information in the model — it leaves a fingerprint carved into the ice.

\section{A constant-free, \ensuremath{\Delta}T-free field test}

A Stefan-number-free dimensional bridge ($\rho L v = q$, $q=-k_{th}(\Delta T/L_{0})\partial \theta /\partial y$, length anchored by the corrugation wavelength) gives $\tau  \sim  \lambda ^{2}/\Delta T$ (amplitude e-fold) and $c_{mig} \sim  \Delta T/\lambda $ (migration speed). In the product the constants $k_{th}$, $\rho _{i}L$ \textbf{and $\Delta T$ cancel exactly}:

\begin{quote}
$I \equiv  \tau \cdot c_{mig}/\lambda  = Im(s)/(2\pi |Re(s)|)$ — the dimensionless ratio of migration rate to amplitude rate, basal-$\Delta T$-free.
\end{quote}

The solver pins $I=O(0.3\mbox{--}0.9)$ ($0.33, 0.68, 0.88$ at $n_{w}=8,12,16$ — wavelength-dependent by a factor $\approx 2.7$, an $O(1)$ morphological number rather than a universal constant).

\textbf{Decoy warning (a methodological result).} $\tau  \propto  \lambda ^{2}/\Delta T$ \emph{looks} like classic curvature/Mullins–Sekerka smoothing, but the $\lambda ^{2}$ comes entirely from the conduction length\ensuremath{\leftrightarrow}time anchoring $L_{0}^{2}/\kappa $, \textbf{not} from a curvature operator — $\beta /a$ is $K$-independent (§3). Fitting $\tau \propto \lambda ^{2}$ in the field would \emph{wrongly confirm} the very curvature ansatz the harmonic probe falsified; the only discriminator is the $K$-exponent of $\beta /a$ ($\approx 0$ here vs $+2$ for curvature), which the dimensional scaling hides. Anyone analysing field scallops should measure the $K$-exponent, not fit $\tau \propto \lambda ^{2}$.

\textbf{Field-test protocol.} On one scallop train measure $\lambda $ (crest spacing), $c_{mig}$ (downstream crest displacement between two visits), and $\tau $ (amplitude e-fold). Then:

\begin{center}\small
\begin{tabularx}{\linewidth}{>{\raggedright\arraybackslash}X>{\raggedright\arraybackslash}X}
\hline
outcome & reading \\
\hline
$I_{obs} \in  [0.2,1.0]$, $c_{mig}$ downstream & consistent with $s=-\beta +i\omega _{mig}$ \\
$c_{mig} \approx  0$ ($Im(s)=0$) & parity-symmetric closure — K-theory \textbf{not} falsified on the ice (contradicts the solver) \\
$c_{mig}$ upstream & new physics (e.g. depositional, not melt-limited) \\
$I_{obs} \gg  1$ or $\ll  0.1$ & $\lambda  \leftrightarrow  u_*$ anchoring mismatch \\
\hline
\end{tabularx}
\end{center}

Because $I$ is $\Delta T$-free, a failure cannot be excused by an unknown basal $\Delta T$; measuring $c_{mig} + \lambda $ alone pins $\Delta T$, so $\tau $ becomes a prediction and the system is over-determined.

\section{Literature grounding}

\begin{itemize}
\item \textbf{$\lambda $ (wavelength) — solid.} $Re_*=u_*\lambda /\nu $ selection is mainstream: Blumberg \& Curl (1974) $\approx 2200$, Thomas (1979) $\approx 1000$, Hsu–Locher–Kennedy (1979) $\lambda =3180 \nu /u_*$, Thorsness–Hanratty (1979)/Hanratty (1981) band $3100\mbox{--}6300 \nu /u_*$.
\item \textbf{$c_{mig}$ direction \& mechanism — solid.} Downstream migration with maximum ablation downstream of the trough at reattachment: Curl (1966), Ashton \& Kennedy (1972), Hsu et al. (1979), Gilpin et al. (1980). Gilpin et al.'s linear stability predicts downstream migration when the heat-flux/topography phase shift lies in $(\pi /2, \pi )$ — the same statement as $Im(s)\ne 0$.
\item \textbf{$c_{mig}$ magnitude and $\tau $ — sparse.} Cave-scallop surveys report $\lambda $ only; migration speed and amplitude decay time are essentially never tabulated together for one train.
\end{itemize}


\textbf{Two derived consistency relations with the turbulence closure (Paper 1).} The migration quadrature $Im(s)$ that distinguishes a real closure is the imaginary part of a single complex transport admittance $Z=a_{r}+ia_{i}$, whose real part is the amplitude-rate / backscatter of the closure (Paper 1; Kraichnan 1976) and whose imaginary part is this scallop migration; a down-gradient eddy diffusivity is the projection $Z\mapsto \max(\mathrm{Re}\,Z,0)$, which zeroes migration and backscatter \emph{together}. Because the flux response is moreover causal, $\mathrm{Re}\,Z(k)$ and $\mathrm{Im}\,Z(k)$ are Kramers--Kronig (Hilbert) conjugates: any \emph{scale-dependent} eddy viscosity --- which every real closure (spectral eddy viscosity, Smagorinsky $\propto |S|$) is --- \textbf{forces} a nonzero migration, so K-theory's pairing of a scale-dependent $\mathrm{Re}\,Z$ with $\mathrm{Im}\,Z\equiv 0$ is acausal (an $O(1)$ Kramers--Kronig residual). The measured migration $I=|\mathrm{Im}\,Z|/(2\pi |\mathrm{Re}\,Z|)$ is therefore not an incidental fact but the causality-mandated shadow of the same operator that carries backscatter in the turbulence closure --- a derived link, verified on the admittance model (discrete Kramers--Kronig reconstruction of $\mathrm{Im}\,Z$ from $\mathrm{Re}\,Z$ to $5.3\%$).

\section{The one open gate — and a regime-matched bound}

\textbf{Sign caveat.} Real ice ripples \emph{grow} ($Re(s)>0$, the Ashton–Kennedy/Gilpin/Hanratty moving-boundary instability), whereas the frozen-interface excess here is decay-only ($Re(s)<0$) by construction. The migration $Im(s)$ is the robustly shared, K-theory-falsifying piece; matching the \emph{growing} branch needs an interface-growth mechanism (a Röthlisberger opening) this smoothing-only solver does not contain.

\textbf{The migration transfers to a moving boundary.} We close the part of this gate the solver \emph{can} reach. Seeding a single mode on the interface and advancing it under the Stefan melt feedback ($\rho L v=q$), the co-evolving interface tracks a \textbf{coherent damped-migrating eigenmode}: $log|Z|$ and $arg(Z)$ of the fundamental shape mode are both linear in time ($R^{2}\ge 0.91$ for amplitude and phase, two seeds \ensuremath{\times} two Stefan numbers). It \textbf{decays} ($Re(s)<0$ at every $St$ — the decay-only branch is \emph{physical}, not an artefact of freezing) and \textbf{migrates downstream}, with $Im(s)$ tracking the frozen-probe quadrature $E_{cos}$ (predicted $Im=E_{cos}/(aSt)\approx -0.43$ vs measured $-0.40$ at $St=10^{-3}$, \textasciitilde{}5–10\%). The observable $I=|Im|/(2\pi |Re|)\approx 0.15\mbox{--}0.17$ is \textbf{$St$-invariant} (the rate scales as \texttt{1/\allowbreak{}St}, so $Re\cdot St$ and $Im\cdot St$ are constant to \textasciitilde{}3\%), confirming $I$ is a well-defined, normalization-free property of the mode. The migration — the closure-distinguishing quadrature — therefore \textbf{survives interface motion}. The total moving-interface decay runs $\sim 1.75\times $ faster than the flow-off conduction estimate (turbulence adds effective smoothing), so this motion-consistent $I\approx 0.16$ sits \emph{below} the frozen flow-off band and \emph{closer to} the Bushuk field bound below. The \emph{growing} branch ($Re(s)>0$) is still out of reach here. (\texttt{scallop\_\allowbreak{}moving\_\allowbreak{}boundary\_\allowbreak{}check.\allowbreak{}py}, \texttt{subglacial/\allowbreak{}moving\_\allowbreak{}boundary\_\allowbreak{}check.\allowbreak{}json}.)

\textbf{A regime-matched real-data bound.} Bushuk et al. (2019, \emph{JFM} 873) track an evolving ice–water interface $h(x,t)$ at sub-mm/15 Hz through transition\ensuremath{\to}equilibrium\ensuremath{\to}\textbf{adjustment}. Their adjustment regime (after the drive was cut $U=1.00\to 0.16 m s^{-1}$) \textbf{damps and migrates downstream simultaneously} — the exact analogue of the frozen probe ($s=-\beta +i\omega _{mig}$, $Re(s)<0$, $Im(s)\ne 0$), so the sign caveat does not apply. Reading their measured kinematics ($c_{mig}=0.11 mm min^{-1}=1.83\times 10^{-6} m s^{-1}$, observed $\lambda \approx 13 cm$, amplitude e-fold $\tau \approx 1\mbox{--}3 h$ by eye):

\begin{quote}
$I_{obs} \approx  0.05\mbox{--}0.15$ (point $\approx 0.10$) at the observed $\lambda $; $O(0.05\mbox{--}0.4)$ overall; \textbf{downstream sign}. Measured the \emph{same way} as the field data — a moving boundary with total (turbulent) smoothing, via the identical $Re(s)=d/dt ln a_{k}$, $Im(s)=d\phi _{k}/dt$ protocol — the solver gives the motion-consistent $I\approx 0.16$ (above), a factor $\approx 1.6$ above the point estimate, versus the factor $\approx 3.3$ against the frozen flow-off band $[0.33,0.88]$. The motion-consistent comparison thus \textbf{reduces the tension to mild} — $I_{obs}$ is $O(0.1\mbox{--}1)$, neither $\approx 0$ nor $\gg 1$, with the predicted downstream sign — a consistency check, not a falsification.
\end{quote}

\textbf{The pin.} The one missing input is Bushuk's underlying $h(x,t)$ arrays, available from the authors on request (\texttt{mitchell.\allowbreak{}bushuk@noaa.\allowbreak{}gov}); the published JFM supplement is a parameter table only, not the interface time series, and no open data deposit exists. The §6 ratio is computed from morphology alone, and the reanalysis harness is \textbf{already implemented and self-validated}: \texttt{harmonic\_\allowbreak{}mode\_\allowbreak{}rate(x,t,H)} rFFTs each frame, tracks the dominant mode $a_{k}(t)e^{i\phi _{k}(t)}$, fits $Re(s)=d/dt ln a_{k}$, $Im(s)=d\phi _{k}/dt$, and returns $I$ plus the downstream sign; on a synthetic damped, downstream-migrating train with \textbf{known} $(Re(s), c_{mig}, \lambda )$ it recovers all three to $<2\%$ and $I$ to $<3\%$. It ingests Bushuk's arrays directly — no further code, just the data. (\texttt{scallop\_\allowbreak{}field\_\allowbreak{}test.\allowbreak{}py}.)

\section{Discussion and scope}

This is a solver-verified, quantitative characterization of the closure: the corrected mode, the operator-swap shortfall, and the temperature-free ratio are all self-consistent and GPU-reproducible. It is \textbf{not yet validated against field scallops} — §8 is precisely why that test is cheap (no $\Delta T$ needed). The banked claims are the qualitative three — $Im(s)\to 0$ at $U=0$, strictly sub-linear migration, exponent $\to \tfrac{1}{2}$ at high $Re_*$ (central value $0.52$) — robust to the \textasciitilde{}±0.1 scatter from spin-up/seed/$U$-grid. The solver-side moving-boundary measurement (§8) is now done — the migration transfers to an actually-moving interface and $I$ is motion-invariant — so what a field dataset would still pin is the migration exponent's precise value and the field $I_{obs}$ from real $h(x,t)$.

\section{Reproduce}

\begin{verbatim}
python glaciers/scallop_amplitude_harmonics.py # §3–§4; writes figures/56_*.json
python glaciers/scallop_ktheory_control.py # §5 keystone control (CPU; xp=cupy for GPU)
python glaciers/scallop_field_test.py # §6/§8; figures/58_*.json (harness self-check + Bushuk bound)
python glaciers/scallop_moving_boundary_check.py # §8 moving-boundary migration transfer + I(St-invariance)
pytest glaciers/tests/test_scallop_amplitude_harmonics.py \
       glaciers/tests/test_scallop_ktheory_control.py \
       glaciers/tests/test_scallop_field_test.py
\end{verbatim}

\codedataavailability{All scallop-migration and harmonic-decomposition results regenerate from the scripts in the reproduce section; the field-test harness and its synthetic validation data are in the public repository at https://github.com/abdh0saifelden2-debug/K.}

\authorcontribution{A. Saifelden conceived the study, performed all computations and analysis, and wrote the manuscript.}

\competinginterests{The author declares that there are no competing interests.}

\begin{thebibliography}{}
\bibitem{ref1}
Ashton, G. D., \& Kennedy, J. F. (1972). Ripples on underside of river ice covers. \emph{Journal of the Hydraulics Division (ASCE)} 98, 1603–1624. doi:10.1061/JYCEAJ.0003407

\bibitem{ref2}
Blumberg, P. N., \& Curl, R. L. (1974). Experimental and theoretical studies of dissolution roughness. \emph{Journal of Fluid Mechanics} 65, 735–751. doi:10.1017/S0022112074001625

\bibitem{ref3}
Bushuk, M., Holland, D. M., Stanton, T. P., Stern, A., \& Gray, C. (2019). Ice scallops: a laboratory investigation of the ice–water interface. \emph{Journal of Fluid Mechanics} 873, 942–976. doi:10.1017/jfm.2019.398

\bibitem{ref4}
Curl, R. L. (1966). Scallops and flutes. \emph{Transactions of the Cave Research Group of Great Britain} 7, 121–160.

\bibitem{ref5}
Gilpin, R. R., Hirata, T., \& Cheng, K. C. (1980). Wave formation and heat transfer at an ice-water interface in the presence of a turbulent flow. \emph{Journal of Fluid Mechanics} 99, 619–640. doi:10.1017/S0022112080000791

\bibitem{ref6}
Hanratty, T. J. (1981). Stability of surfaces that are dissolving or being formed by convective diffusion. \emph{Annual Review of Fluid Mechanics} 13, 231–252. doi:10.1146/annurev.fl.13.010181.001311

\bibitem{ref7}
Hsu, K.-H., Locher, F. A., \& Kennedy, J. F. (1979). Forced-convection heat transfer from irregular melting wavy boundaries. \emph{Journal of Heat Transfer} 101, 598–602. doi:10.1115/1.3451043

\bibitem{ref8}
Mullins, W. W., \& Sekerka, R. F. (1964). Stability of a planar interface during solidification of a dilute binary alloy. \emph{Journal of Applied Physics} 35, 444–451. doi:10.1063/1.1713333

\bibitem{ref9}
Thomas, R. M. (1979). Size of scallops and ripples formed by flowing water. \emph{Nature} 277, 281–283. doi:10.1038/277281a0

\bibitem{ref10}
Thorsness, C. B., \& Hanratty, T. J. (1979). Mass transfer between a flowing fluid and a solid wavy surface. \emph{AIChE Journal} 25, 686–697. doi:10.1002/aic.690250415

\end{thebibliography}

\end{document}
